\section{From Interface Error to Return Error}
\label{sec:theory}

Identification does not quantify finite-sample error. For a finite horizon \(H\), compare the true primitive \(K\) with \(\widehat K\) while querying the same \(\pi\) and \(K_\Pi\). At each \((x,u)\), suppose a measurable one-step coupling of \((x',c)\) and \((\widehat x',\widehat c)\) exists, and learned continuation laws are \(L_{t+1}\)-Lipschitz in \(W_1\). Define
\[
\epsilon_t(x,u)=\mathbb E\big[|c-\widehat c|+L_{t+1}d(x',\widehat x')\big].
\]

\begin{theorem}[Occupancy-weighted finite-horizon transport]
\label{thm:transport}
Under the coupling and continuation regularity above,
\begin{equation}
W_1(\mathcal Z_{0:H}(x_0),\widehat{\mathcal Z}_{0:H}(x_0))
\leq\sum_{t=0}^{H-1}\mathbb E_{(x_t,u_t)\sim d_t^{\pi,\Pi,K}}[\epsilon_t(x_t,u_t)].
\label{eq:transport}
\end{equation}
Moreover, the proper SSP gives
\begin{equation}
W_1(\mathcal Z_C(x),\mathcal Z_{C,H}(x))
\leq c_{\max}\mathbb E[(T_C-H)_+]\to0.
\label{eq:truncation}
\end{equation}
\end{theorem}

Coverage alone does not imply Equation~\eqref{eq:transport}: continuation regularity may fail near discontinuous safety-filter active sets. The bound instead motivates two empirical variables---executed-interface extrapolation and hitting horizon---and requires both to be controlled in the experiment.

Finally, let the oracle and learned threshold decisions be
\(d^\star=\mathbf{1}\{b\leq U+m\}\) and
\(\widehat d=\mathbf{1}\{b\leq\widehat U+m\}\).

\begin{theorem}[Return-boundary stability]
\label{thm:boundary}
If \(|\widehat U-U|\leq\varepsilon\), then
\begin{equation}
d^\star\neq\widehat d\quad\Longrightarrow\quad |b-m-U|\leq\varepsilon.
\label{eq:boundary}
\end{equation}
\end{theorem}

Thus error can alter the threshold decision only within an \(\varepsilon\)-band around the boundary. This does not make global MAE irrelevant, nor imply that every disagreement causes stranding; it identifies where underestimation becomes decision-critical. Appendix~\ref{app:proofs} includes the proofs, decision-check overshoot condition, and task-plus-return risk allocation.
